 \providecommand{\main}{../../..}
\documentclass[\main/main.tex]{subfiles}
\begin{document}
\section{Teorema di dualità forte}

\begin{minipage}{\textwidth}
  \begin{minipage}{.83\textwidth}
    \flushleft
    \begin{theorem}[Dualità forte]
      \label{dualita_forte}
      Sia $P = \crl{\bmx \geq \bm{0}: A\bmx \geq \bmb} \neq \emptyset$ con soluzione ottima \textbf{finita}. Allora la soluzione ottima del primale coincide con la soluzione ottima del duale.

      \begin{figure}
        \[
          \min\crl{\bmct\bmx: A\bmx \geq \bmb, \bmx \geq \bm{0}} = \max\crl{\bmut\bmb: \bmct \geq \bmut A, \bmu \geq \bm{0}}
        \]
        \caption{Teorema di dualità forte: la soluzione ottima primale coincide con la soluzione ottima duale}
      \end{figure}
    \end{theorem}
  \end{minipage}\hfill
  \begin{minipage}{0.15\textwidth}\center
    \includegraphics[width=\textwidth]{random/salt}
  \end{minipage}
\end{minipage}

\begin{proof}
  Segue direttamente dal lemma di Farkas (\ref{lemma_farkas}).
\end{proof}
\end{document}